\documentclass[11pt]{article}
\usepackage[margin=1in]{geometry}
\usepackage{amsmath,amssymb}
\usepackage{enumitem}
\usepackage{booktabs}
\usepackage[hidelinks]{hyperref}
\usepackage[round]{natbib}
\setlength{\parskip}{0.5em}
\setlength{\parindent}{0pt}

\newcommand{\rev}[1]{\par\medskip\noindent\textbf{#1}\par}
\newenvironment{comment}{\par\medskip\begin{quote}\small\itshape}{\end{quote}}
\newcommand{\reply}{\noindent\textbf{Response.}\ }

\begin{document}

\begin{center}
{\large\bfseries Response to the reviewers}\\[0.3em]
Manuscript: \emph{Minimum Weighted Hazard-Exposure Dispatch: A Scheduling
View with Verified Algorithms and a Wildfire Case Study} (formerly
\emph{Minimum Weighted Hazard-Exposure Dispatch: Complexity, an Exact
Algorithm, and an FPTAS})\\
Journal of Combinatorial Optimization --- major revision
\end{center}

Dear Editor and Reviewers,

Thank you for the careful reading and the constructive reports. We
revised the manuscript substantially. Reviewer~1's central objection,
that the theoretical results follow from the classical theory of
$1\|\sum w_jU_j$, is correct, and we have rewritten the title, abstract,
introduction, related work and summary table to say so plainly, to cite
the classical source for each result, and to separate what is classical
from the few things we believe are new. Reviewer~1's question about the
FPTAS exposed a real weakness in our experiments (the FPTAS's rounding was
inactive in most of them), which we explain and fix. Reviewer~2's
requests for a clearer hardness proof, a self-contained proof of the
equal-cost theorem, a shorter positioning section and a discussion of
parameter uncertainty are all addressed.

Before resubmitting we re-checked every numerical result, every cited
source, and the proofs (the statements and proofs are now also
machine-checked in Lean~4), and we had the revision critically reviewed
internally. This found several errors of our own, some in the first
version and some introduced while revising; we list all of them in
Section~\ref{sec:own}. The most consequential are: the case study had
used extrapolated hazard times that the cited NIST report in fact
documents (and that were far off), and had a modelling premise that its
own geography contradicts; an extension we had presented as new (the
equal-cost theorem for several vehicles) is classical; and our
``tight'' FPTAS guarantee was not the best provable one. All tables and
figures were regenerated.

\paragraph{Numbering.} Items are numbered in order of appearance, so the
numbers of the submitted version have shifted. Submitted $\to$ now:
Lemma~1 $\to$ Lemma~2; Theorem~2 (hardness) $\to$ Theorem~3; Theorem~3
(dynamic program) $\to$ Theorem~5; Theorem~4 (FPTAS) $\to$ Theorem~7;
Theorem~5 (equal cost) $\to$ Theorem~8; Propositions~6, 7, 8 $\to$
Propositions~9, 10, 11; Theorem~9 (classification) $\to$ Theorem~12.
New: Proposition~1 (the round-trip model is conservative),
Proposition~4 (constant hazard-arrival time), Lemma~6 (value-indexed
table). Fact~1 keeps its number. Below, ``Theorem~$a$ ($b$)'' means
Theorem~$a$ now and Theorem~$b$ in the submitted version. Section
numbers: Section~5 now has a new Subsection~5.5 (general-purpose solver),
so robustness is 5.6 and the Camp Fire study is 5.7.

\section*{Summary of the main changes}

\begin{enumerate}[leftmargin=*,itemsep=2pt]
\item \textbf{Positioning and novelty} (title, abstract, Sections~1, 2,
Table~2, conclusion): every result that is classical is labelled as such
with its source (Karp 1972; Lawler and Moore 1969; Sahni 1976; Gens and
Levner 1981; Lawler 1976; Dessouky et al.\ 1990; Moore 1968), and the
claims to novelty are cut down to what we can support.
\item \textbf{What we now claim} (Section~4 and Table~2): a refined,
exactly tight analysis of the FPTAS (Theorem~7, Proposition~11); hardness
under the arrival reading even when the hazard-arrival time is constant
(Proposition~4); the observation that the depot round-trip model is
conservative with respect to a chained route (Proposition~1); and
elementary worst-case examples for natural heuristics (Propositions~9
and~10). The equal-cost theorem (Theorem~8) now covers vehicles with
different speeds and is labelled classical; the classification
(Theorem~12) is labelled an elementary corollary.
\item \textbf{Hardness proof} (Theorem~3 (2)): a direct reduction from
\textsc{Partition} with an explicit lemma (Fact~1) on feasibility under
a common deadline.
\item \textbf{Experiments} (Section~5): redesigned, with exhaustive-search
verification, instance sizes to $n=400$, a stronger baseline, FPTAS
experiments in which rounding is active, an adversarial family, and a
comparison with a general-purpose MIP solver.
\item \textbf{Case study} (Section~5.7): rebuilt on hazard-arrival times
documented by NIST for every community (nothing is extrapolated), with a
comparison against a chained-route model, a sensitivity analysis with a
stated protocol, and a sample-average plan.
\item \textbf{Discussion} (Section~6): corrected multi-vehicle
subsection; parameter uncertainty in a new subsection with the two
suggested references.
\item \textbf{Reproducibility}: all code and a Lean~4 formalisation of
the theory are public (Data Availability).
\end{enumerate}

\section{Reviewer 1}

\begin{comment}
While the manuscript is technically correct and generally well
presented, the level of novelty of the theoretical contributions is not
sufficiently convincing. [\dots] most of the main theoretical results
appear to follow from well-known results in the scheduling literature.
[\dots] Consequently, the manuscript does not clearly establish new
complexity results, new algorithmic techniques, or substantially new
structural properties beyond those already known for the equivalent
scheduling formulation.
\end{comment}

\reply We agree. The manuscript proves that MWHED is
$1\|\sum w_jU_j$ and the general results follow from it. We have not
tried to argue otherwise. Concretely:

\begin{itemize}[leftmargin=*,itemsep=2pt]
\item The title, abstract and Section~1 now say that the basic
complexity picture and the algorithms are inherited from scheduling
theory and that we claim no new technique for the general problem. The
contributions are split into a list of classical results restated with
complete proofs (hardness, the dynamic program, the FPTAS, the
equal-size greedy, the classification) and a list of modest additions.
\item Section~2 states, for each result, which classical result it is:
Theorem~5 (3) is the Lawler--Moore dynamic program; Theorem~7 (4) is the
classical value-scaling scheme for job sequencing with deadlines
(\citealt{sahni1976}, \citealt{gensl1981}); the hardness is Karp's
\citep{karp1972} and the knapsack equivalence; and Theorem~8 (5) is the
scheduling of unit-time jobs on identical or uniform machines, for which
$O(n\log n)$ algorithms are known \citep{lawler1976,dessouky1990}.
Our submitted version presented the extension of Theorem~8 to several
vehicles as new; it is not, and we now say so (see Section~\ref{sec:own}).
\item The summary table (Table~2) has a column giving the origin of each
result (classical, Lawler--Moore, Moore, elementary corollary, new).
\item We removed the claim that the dynamic program and the FPTAS are
``necessary'' for any guarantee: Propositions~9 and~10 show only that
deadline-order dispatch and one greedy repair have no guarantee.
\end{itemize}

\begin{comment}
1. Could the authors identify a new structural property?
\end{comment}

\reply Yes, but they are modest, and we present them as such.

\begin{itemize}[leftmargin=*,itemsep=2pt]
\item \textbf{Proposition~4 (constant hazard-arrival time).} Under the
arrival reading of Remark~1 the natural homogeneous case is a hazard that
reaches all sites at the same time, which is \emph{not} a constant
deadline $d_i$ (it is $d_i=H+p_i/2$). We prove that MWHED stays weakly
NP-hard there, even with all $p_i$ even and $w_i=p_i/2$; constant $p$ or
constant $w$ remain polynomial, so the restriction cannot be sharpened.
\item \textbf{Proposition~1 (the round-trip model is conservative).} If
the sites lie in a metric space, a crew that chains the sites reaches each
site no later than in the depot round-trip model, so the round-trip
optimum is a lower bound on what a crew can protect. This makes explicit
when the premise of the model is costly (clustered sites) and is
quantified on our case study.
\item \textbf{A sharper, exactly tight FPTAS analysis} (Theorem~7,
Proposition~11): the scheme of Theorem~7 guarantees a constant
$\rho_{n,\epsilon}$ that exceeds both $1-\epsilon$ and $1/(1+\epsilon)$,
and the constant is attained (two extremal families), so the classical
$1-\epsilon$ bound is not attained.
\item \textbf{The classification} (Theorem~12) remains, but we now call it
an elementary corollary, and we note that the natural tractability
boundary in the scheduling literature is Lawler's agreeable-weights case
\citep{lawler1976}, which contains both constant $p$ and constant $w$.
\end{itemize}

\begin{comment}
2. The manuscript repeatedly states the dynamic program and FPTAS as
major contributions. However, the relationship between these algorithms
and classical Lawler-Moore-type methods is not discussed in sufficient
detail. Are the algorithmic contributions truly different from known
scheduling algorithms?
\end{comment}

\reply No. Theorem~5's dynamic program is the Lawler--Moore dynamic
program (jobs in due-date order, state equal to the accepted processing
time, one test per accepted job); Remark~2 says so and describes exactly
what our proof adds (a derivation from the earliest-deadline-first lemma
in the notation of the dispatch problem). Theorem~7 is the classical
value-scaling scheme of Sahni (1976) and Gens and Levner (1981); the
manuscript states this, notes that faster schemes are known, and keeps
the classical construction because its analysis is the cleanest and,
with Theorem~7 and Proposition~11, exactly tight. Neither algorithm is
presented as a contribution; they are a self-contained treatment for a
transportation audience.

\begin{comment}
3. Why does FPTAS = 1000 always appears in every table? Could the
authors strengthen the empirical evaluation?
\end{comment}

\reply This was a genuine weakness. With weights at most $100$ the
scaling factor $K=\max(1,\epsilon w_{\max}/n)$ equals $1$ in almost all
of our runs, so no rounding took place and the ``FPTAS'' coincided with
the exact value-indexed dynamic program. A ratio of $1.000$ was therefore
guaranteed and said nothing about the scheme. We state this explicitly
(Section~5.2, last row of Table~4: the share of instances with $K=1$;
here $n$ is the number of sites after deleting individually infeasible
ones) and have redesigned the study (every table was regenerated):

\begin{itemize}[leftmargin=*,itemsep=2pt]
\item \textbf{Verification} (Section~5.1): every algorithm is checked
against exhaustive search, including all $n!$ dispatch orders for
$n\le7$ (400 instances), best-feasible-subset enumeration for $n\le14$
(1,500 instances, one third with weights up to $10^6$, FPTAS guarantee
checked in $4{,}500$ runs of which $2{,}923$ had active scaling), and
exhaustive search over assignments for the equal-cost algorithm with
$m=2,3$ vehicles; the added results are checked by a further script
(1,600 instances for the different-speed theorem, 600 for
Proposition~4, 4,400 random instances plus an adversary for the refined
FPTAS bound, and recursions cell by cell). No violations.
\item \textbf{Accuracy across sizes} (Section~5.2): $n$ up to $400$, $50$
instances per size, worst-case ratios as well as means, and a stronger
baseline, EDD with skipping, in addition to naive EDD.
\item \textbf{FPTAS with active scaling} (Section~5.3): weights up to
$10^6$ ($K$ in the hundreds to thousands), a strongly correlated family,
and an adversarial family. On random and strongly correlated instances
the FPTAS is nearly exact (worst ratio $0.9993$ at $\epsilon=0.5$)
because the bound is pessimistic when the optimum contains many sites; on
the adversarial family the ratio falls to $0.669$ at $\epsilon=0.5$ and
follows the guarantee. We now state that the adversarial column is a
single deterministic instance equal to a closed form, not an empirical
worst case. A simple completion step repairs the family completely.
\item \textbf{Runtime} (Section~5.4) re-run with active scaling, and a
\textbf{comparison with a general-purpose solver} (Section~5.5): HiGHS on
the subset formulation returns the same optimum on all $40$ instances and
is 15 to 200 times slower than the dynamic program.
\end{itemize}

\section{Reviewer 2}

\begin{comment}
3. The part of MWHED's positioning in page 6 is comprehensible,
nevertheless it could be rephrased better by shorter arguments.
\end{comment}

\reply The ``Positioning'' paragraph (Section~2) was shortened and now
lists the literatures it draws on and refers to Table~2; the two long
paragraphs on real-time schedulability and matroid-partition algorithms
were cut to two sentences and removed, respectively, since they carried
no result used later.

\begin{comment}
1. Theorem 2 on NP-hardness of MWHED might require a bit efforts to get
full clarity. In the 3rd line from the bottom of page 9, don't you think
that the fact ``completion times along a fixed set are additive'' is too
strong? What fixed sets should we use here? [\dots] The authors should
explain about additivity of special sets [\dots] since it is not obvious
by default, feasibility condition does not explicitly implies
additivity. [\dots] the above condition ($\diamondsuit$)
$\sum_{i\in S}p_i\le D$ is crucial to relate your formulation to a 0/1
knapsack instance.
\end{comment}

\reply Thank you; the sentence was imprecise. We replaced it by an
explicit statement, \textbf{Fact~1} (Section~4.1), with a proof:
\emph{if all deadlines equal $D$, a set $S$ of sites can all be served on
time iff $\sum_{i\in S}p_i\le D$}, which is exactly your condition
($\diamondsuit$). The ``fixed set'' is an arbitrary subset $S$. The only
additivity used is that the \emph{last} site served from $S$ completes
at the sum of all of $S$'s occupancies, because the vehicle is never idle
and each $p_i$ does not depend on the predecessor; this gives necessity.
For sufficiency every completion time is a partial sum of positive
numbers and so is at most the total. The proof of Theorem~3 cites Fact~1,
and the knapsack correspondence for general weights is stated after the
theorem. We also made the reduction handle odd total $A$ and sites with
$a_i>A/2$ explicitly.

\begin{comment}
2. [\dots] your claim on ``NP-complete of the 0/1 knapsack problem''
comes from polynomial reduction of PARTITION (one of the six popular
graph-based NP-complete problems [\dots]). This is a solidly smart
argument, if the above concern is clearly elucidated. Otherwise, a
graph-based pure proof for Theorem 2 possibly might give better insights
(ref. Delazeri and Ritt (2026)).
\end{comment}

\reply We removed the dependence on the knapsack citation: Theorem~3
gives a \emph{direct} polynomial reduction from \textsc{Partition} to
MWHED ($p_i=w_i=a_i$, $d_i=A/2$), so that the proof is self-contained
and shows hardness when $w_i=p_i$; membership in NP is added. A small
remark on terminology: \textsc{Partition} is one of the basic
NP-complete problems in Garey and Johnson, but it is a number problem,
not a graph problem. We did not give a graph-based proof: MWHED has no
graph structure once each visit is a round trip from the depot, so there
is no natural graph on which to base one.
\citet{delazeri2026} prove NP-completeness for a different, graph-based
model of wildfire suppression with simultaneous resources; Section~2 says
so.

\begin{comment}
3. [\dots] the authors would elucidate more explicitly Theorem 5's
proof. In particular, would you remind readers the concepts (a)
transversal matroid, union-find structure, and (b) the related notation
$\alpha(n)$?
\end{comment}

\reply Section~4.4 opens with a paragraph recalling the notions used:
matroid and its three axioms, the greedy algorithm, the transversal
matroid (sites against (vehicle, position) slots), and the union--find
structure, with $\alpha$ defined as the inverse Ackermann function. The
proof of Theorem~8 (5) is in five steps: (1) slots; (2) a counting
criterion for feasibility; (3) a direct proof of the matroid augmentation
property; (4) a short proof that greedy is optimal on any matroid; (5) the
implementation, with a proof that it accepts a site exactly when the kept
set plus that site is feasible. The theorem now allows vehicles with
different dispatch times and is stated as classical. In answering this
comment we also corrected our union--find remark: the
$O(n\alpha(n))$ bound needs union by size or rank in addition to path
compression (path compression alone gives $O(\log n)$ amortized), and we
describe the label array that gives it for our linking direction.
Algorithm~3 and Example~5 are rewritten accordingly.

\begin{comment}
4. [\dots] do we agree that the problem of MWHED is obviously
data-driven, the proofs and optimal solutions obtained depend on both
accurate formulation and precise calibration of model parameters, upon
uncertainty? To some extent, would formulation of uncertainty [\dots]
(at not just variable level but also at parameter one) be considered in
next studies, where stochastic programming might be a potential
approach? Kindly, e.g. refer to problems of evacuation planning under
unsafe contexts and facility location under disruptive events (FLDE),
ref. [1, 2].
\end{comment}

\reply We agree. Section~6.3 (``Uncertain parameters'') covers
uncertainty in $p_i$, $d_i$ and $w_i$, not only in the deadlines. It
discusses (i) scenario-based and two-stage stochastic formulations and
which of our results survive (the earliest-deadline lemma holds scenario
by scenario but not simultaneously; a worked example shows that sorting
by expected deadline can be nearly twice as bad as optimal), citing your
two references \citep{wang2020evac,koca2023}; (ii) a robust-feasibility
observation; and (iii) calibration and sensitivity analysis. The case
study (Section~5.7) now includes a Monte-Carlo sensitivity analysis over
response speed, documented event times and departure delay, in an
independent and in a correlated variant, with standard errors: the plan
that is optimal for the nominal instance is fully on time in only
$16.2\%$ of $20{,}000$ draws and retains $31.9\%$ of the drawn optimum on
average, whereas a plan chosen against a sample of perturbed scenarios is
fully on time in $79.3\%$ and retains $99.7\%$. Stochastic programming
for MWHED itself is listed as the main direction for future work; we have
not attempted a solution method for it here.

\section{Errors of our own corrected in this revision}\label{sec:own}

In the course of revising we re-checked the manuscript, the code, the
cited sources and the proofs, and found the following, which we correct
and disclose.

\begin{enumerate}[leftmargin=*,itemsep=3pt]
\item \textbf{Case study: hazard times.} Two of the four hazard-arrival
times had been \emph{extrapolated} by an isotropic spread-rate rule
(Magalia at $58$ and Yankee Hill at $64$ minutes), although the cited
NIST report documents events for them: spot fires in Old Magalia at 08:40
and the fire running through Yankee Hill on the following day. The
extrapolation was far off. We also had attributed a ``time zero'' of the
timeline to NIST, but the report uses clock times only; the earlier
anchors of $52$ and $71$ minutes were reproducible only with a start time
taken from another source. The case study now uses the documented clock
times of the first events in the four communities, with time zero equal
to the initial dispatch time in NIST's dispatch log, states the event
definitions as NIST words them, and no longer attributes our ignition
coordinate to NIST (the figure uses NIST's origin point). All case-study
results were recomputed; the earlier numbers ($26{,}218$, $26{,}928$) are
superseded.
\item \textbf{Case study: modelling premise.} Concow, Paradise and Yankee
Hill are within about $10$~km of each other but $22$--$37$~km from the
depot, so the round-trip model of the paper was far from the geometry,
and the earlier conclusion that only Paradise could be protected was an
artefact of that premise. The paper now proves that the round-trip model
is conservative (Proposition~1), compares it with a chained route on the
same data (Table~12), and says that the case study illustrates the model
and does not support operational recommendations.
\item \textbf{Case study: protocol.} The ``retained share'' of the nominal
plan was not reproducible from the text (an independent replication under
another reading of the plan gave about $22\%$ instead of $26.6\%$). The
protocol (fixed visiting sequence, skipping sites deleted by
Assumption~1) is now stated, the alternative reading is reported, standard
errors are given, a correlated-error variant is added, and the weight
at which Paradise would leave the optimum is reported. We also added a
paragraph on the ethics of conceding low-weight sites and on what
``protected'' means.
\item \textbf{Case study: deadline semantics} (from the first revision).
The model's on-time test counts the completion of the round trip (return
reading), but the narrative speaks of reaching the site (arrival
reading). Remark~1 makes the two readings precise (the arrival reading is
an instance of MWHED with deadlines $d_i^{\mathrm{haz}}+p_i/2$) and the
case study uses it; the return reading is reported for comparison.
Individually infeasible sites are deleted before every algorithm,
including the naive baseline; statements that the optimum ``saves'' a
number of people were replaced by the on-time weight.
\item \textbf{Novelty of Theorem~8 (5).} We had presented the extension of
the equal-cost greedy to several identical vehicles as new. It is
classical (Lawler 1976 for identical machines; Dessouky, Lageweg, Lenstra
and van de Velde 1990 for uniform machines); we now cite both, state the
theorem for vehicles with different dispatch times, and keep our proof
for self-containedness. The classification (Theorem~12) is labelled an
elementary corollary.
\item \textbf{Release dates.} The statements that equal-cost tractability
``is not robust to release dates'' and that hardness follows ``once
release dates are added'' were not supported for one vehicle: the cited
result of \citet{heegermolter2025} concerns the unweighted problem on
several identical machines, and for one machine equal processing times
with release dates and weights are polynomial \citep{baptiste1999}. The
text now says exactly this. The earlier inference that the equal-cost
case becomes hard with several vehicles was also wrong, since that
result requires release dates, which MWHED does not have.
\item \textbf{Tightness of the FPTAS.} The submitted text claimed that
the guarantee $1-\epsilon$ is ``tight up to a factor $1-\epsilon^2$''.
The scheme in fact guarantees a larger constant $\rho_{n,\epsilon}$
(Theorem~7), which is attained by two families (Proposition~11); the
classical $1-\epsilon$ is not attained. The statements in the abstract,
introduction, Table~2 and conclusion were rewritten.
\item \textbf{Proofs.} The main-text proof of the FPTAS guarantee used an
inequality in the wrong direction (the appendix repaired it); the displayed
recursions for the dynamic programs, written as an indicator times a
value, were wrong as literally stated (they return $0$ when the test
fails), and are now written with cases; the proof of Lemma~2 did not
cover ties between deadlines; the union--find cost claim needed union by
size (see Reviewer~2, point~3); the reduction of Theorem~3 did not handle
odd $A$ explicitly. Each fix was checked against brute force
(\texttt{verify\_extensions.py}).
\item \textbf{Heeger--Hermelin.} The submitted text said that a small
number of distinct dispatch costs ``does not guarantee tractability''. The
result is $W[1]$-hardness in that parameter (no $f(k)\mathrm{poly}(n)$
algorithm); the text now says this and makes no claim about two or more
distinct values.
\item \textbf{Multiple vehicles.} The claim that the $m=1$ instance
``embeds'' by adding idle vehicles was wrong (extra vehicles can serve
sites); we now use blocker sites.
\item \textbf{Implementation.} The FPTAS guarantee relies on
Assumption~1 (so that $w_{\max}\le W^\ast$), which the code did not
enforce; both algorithms now delete infeasible sites first, and the code
was rewritten and checked against exhaustive search.
\item \textbf{Bibliography.} We verified every entry against a
publisher, Crossref or arXiv record and corrected: the author lists and
volume of \citet{alvelos2025} (three authors, volume 86), the surname
\emph{von der Br\"uggen} (G\"unzel et al.), the author order of the NIST
report, the article type of the AAMAS 2026 entry (extended abstract), and
missing volume, article and page data for several proceedings entries.
\item \textbf{Verification claim.} The Lean~4 statement in our first draft
of the Data Availability section was inconsistent about what is
formalised (complexity-class and running-time statements are not); it now
says precisely what is and is not formalised, and the library, the code
and the experiments are public.
\item Minor: the dash convention in the dynamic-program example table was
described incorrectly; the explanation of why the FPTAS runtime was flat
in $\epsilon$ cited saturating weights, whereas the cause was $K=1$. Both
are fixed.
\end{enumerate}

We believe the revised manuscript is considerably more honest about what
is and is not new, and we thank the reviewers for the comments that led
to these changes.

\bigskip
\noindent Sincerely,\\
The authors

\bibliographystyle{plainnat}
\begin{thebibliography}{12}
\bibitem[Alvelos et al.(2025)]{alvelos2025} F.~Alvelos, M.~Marto, A.~Mendes. Dispatching, positioning and routing resources for wildfire initial attack. \emph{Networks} 86(1):80--102, 2025.
\bibitem[Baptiste(1999)]{baptiste1999} P.~Baptiste. Polynomial time algorithms for minimizing the weighted number of late jobs on a single machine with equal processing times. \emph{Journal of Scheduling} 2:245--252, 1999.
\bibitem[Delazeri and Ritt(2026)]{delazeri2026} G.~Delazeri, M.~Ritt. Wildfire Suppression: Complexity, Models, and Instances. arXiv:2603.29865, 2026.
\bibitem[Dessouky et al.(1990)]{dessouky1990} M.~I. Dessouky, B.~J. Lageweg, J.~K. Lenstra, S.~L. van de Velde. Scheduling identical jobs on uniform parallel machines. \emph{Statistica Neerlandica} 44(3):115--123, 1990.
\bibitem[Gens and Levner(1981)]{gensl1981} G.~V. Gens, E.~V. Levner. Fast approximation algorithm for job sequencing with deadlines. \emph{Discrete Applied Mathematics} 3:313--318, 1981.
\bibitem[Heeger and Hermelin(2024)]{heegerhermelin2024} K.~Heeger, D.~Hermelin. Minimizing the weighted number of tardy jobs is W[1]-hard. ESA 2024.
\bibitem[Heeger and Molter(2025)]{heegermolter2025} K.~Heeger, H.~Molter. Minimizing the number of tardy jobs with uniform processing times on parallel machines. STACS 2025.
\bibitem[Karp(1972)]{karp1972} R.~M. Karp. Reducibility among combinatorial problems. In \emph{Complexity of Computer Computations}, Plenum, 1972.
\bibitem[Koca(2023)]{koca2023} E.~Koca. Two-stage stochastic facility location problem with disruptions and restricted shortages. \emph{Computers \& Industrial Engineering} 183:109484, 2023.
\bibitem[Lawler(1976)]{lawler1976} E.~L. Lawler. Sequencing to minimize the weighted number of tardy jobs. \emph{RAIRO Recherche Op\'erationnelle} 10(5 Suppl.):27--33, 1976.
\bibitem[Sahni(1976)]{sahni1976} S.~K. Sahni. Algorithms for scheduling independent tasks. \emph{Journal of the ACM} 23(1):116--127, 1976.
\bibitem[Wang(2020)]{wang2020evac} L.~Wang. A two-stage stochastic programming framework for evacuation planning in disaster responses. \emph{Computers \& Industrial Engineering} 145:106458, 2020.
\end{thebibliography}

\end{document}
